Shelfmark 148, scan 12.

\scanpage{12}

\begin{center}
U N I V E R S I T Y\\
OF SCIENCE AND TECHNOLOGY\\
OF LANGUEDOC\\
\textasciicircum\textasciicircum\textasciicircum\textasciicircum\textasciicircum\textasciicircum\textasciicircum\textasciicircum\textasciicircum\textasciicircum\textasciicircum\textasciicircum\textasciicircum
\end{center}

Montpellier, 14 October 1983

Registrar's Office

Distribution : List A\\
MA/gs \hfill Number 4470

\begin{center}
NOTICE OF DEFENCE OF A DOCTOR-ENGINEER THESIS

\uuline{`` MATERIALS PHYSICS ''}
\end{center}

\uline{SURNAME and First name of the candidate} : GOUDEAU Jacques

\uline{Title of the thesis} : ``Thermal quenching stresses in parts with rotational symmetry''

\uline{Scheduled date, time and place of the defence} : On 28 October 1983, at 11 o'clock\\
Room of the Institute for Research on the Teaching of Mathematics - University of Science and Technology of Languedoc

\uline{Name of the thesis supervisor} : Mister ZARZYCKI - Mister MAISONNEUVE

\uline{Designation of the Laboratory where the thesis is being prepared} : Laboratory of Materials Science, Continuum Mechanics, Structural Chemistry - University of Science and Technology of Languedoc.

\uline{Analysis of the work} : We study the evolution of the temperature and stress fields, in a glass part with rotational symmetry.

For this we solve by finite elements a problem of thermovisco-elasticity, with a constitutive law of MAXWELL type with coefficients varying as a function of temperature.

We compare the results obtained with those of an optical method of stress analysis as well as with an industrial method of measurement.

\uline{Approximate duration of the research that led to the thesis} : 2 years

\uline{Composition of the proposed examining committee} : \begin{tabular}[t]{ll}
Misters & ZARZYCKI J.\\
& MAISONNEUVE O.\\
& REGNIER J\\
& FRANCOIS\\
& BOURRET D.
\end{tabular}

\begin{flushright}
For the Secretary General,\\
the Principal Officer for\\
School and University Administration,

\unsure{A. Gentil}

\uline{A. GENTIL}
\end{flushright}

\margin{11}
